% Generated draft: requires mathematical and source review.
\documentclass[11pt]{article}
\usepackage[margin=25mm]{geometry}
\usepackage{fontspec}
\setmainfont{Latin Modern Roman}
\usepackage{amsmath,amssymb,mathtools,unicode-math}
\setmathfont{Latin Modern Math}
\usepackage{enumitem,needspace,xurl,hyperref}
\setlength{\parindent}{0pt}
\setlength{\parskip}{5pt}
\setlength{\emergencystretch}{3em}
\newcommand{\N}{\mathbb{N}}
\newcommand{\Z}{\mathbb{Z}}
\newcommand{\Q}{\mathbb{Q}}
\newcommand{\R}{\mathbb{R}}
\newcommand{\C}{\mathbb{C}}
\newcommand{\F}{\mathbb{F}}
\newcommand{\A}{\mathbb{A}}
\newcommand{\PP}{\mathbb{P}}
\newcommand{\E}{\mathbb{E}}
\newcommand{\eps}{\varepsilon}
\newcommand{\dd}{\,\mathrm{d}}
\begin{document}
\section*{UnsolvedMath ID: 30004760}
\textbf{Source problem code:} catalogue entry
\subsection{Definitions and mathematical statement}
The supplied question concerns an extension of a regularity theorem for rectifiable varifolds to anisotropic energies. To interpret the displayed energy in conventional Euclidean notation, let \(\Gamma\) be a countably \(m\)-rectifiable set in an ambient Euclidean space \(\R^n\), let \(T_z\Gamma\) denote its approximate tangent \(m\)-plane at \(\mathcal H^m\)-almost every point, and let \(\theta\) be a nonnegative measurable multiplicity giving a locally finite weight measure
\[
\mu_V(A)=\int_{A\cap\Gamma}\theta(z)\,\mathrm d\mathcal H^m(z).
\]
Here \(V\) is the rectifiable varifold represented by \((\Gamma,\theta)\). Write \(G(m,n)\) for the Grassmannian of unoriented \(m\)-dimensional linear subspaces of \(\R^n\). An integrand \(\Psi\) on tangent planes gives the energy displayed in the source:
\[
\Sigma_\Psi(V)=\int_\Gamma\Psi(T_z\Gamma)\theta(z)\,\mathrm d\mathcal H^m(z),
\]
whenever this integral is defined. These conventions explain the notation; the supplied record does not specify an ambient domain, codimension, integrality of the multiplicity, or an admissible class of integrands.

\emph{Question as supplied.} Does the regularity theorem for rectifiable \(m\)-varifolds whose generalized mean curvature belongs to \(L^p\), with \(p>m\), extend to the anisotropic energy \(\Sigma_\Psi\), under the corresponding density-ratio hypothesis?

\emph{Unresolved qualifications.} The record does not identify the precise isotropic regularity theorem, its regularity conclusion, the density-ratio inequality or its constants, or the quantitative curvature hypothesis. It also does not specify whether curvature integrability is local or global. For the anisotropic extension, the required smoothness, positivity, ellipticity, or other structural conditions on \(\Psi\), and the definition of the curvature quantity controlled by anisotropic first variation, are absent. The record does not say whether singular first variation is excluded. It supplies no interior or boundary formulation. Consequently, this question is a schematic target and cannot be converted into a fully quantified regularity assertion from the supplied record alone.
\subsection{Short English statement}
The question asks whether regularity for rectifiable varifolds with generalized mean curvature in \(L^p\), where \(p\) exceeds the varifold dimension, extends to energies depending on tangent-plane direction. The density-ratio assumption, admissible anisotropies, curvature definition, and intended regularity conclusion require clarification.
\subsection{Sources}
\begin{itemize}
\item Supplied local UnsolvedMath registry record, permanent ID \(30004760\), titled \emph{Regularity of Varifolds for Anisotropic Energies}; source problem code \texttt{catalogue entry}, in the Oberwolfach Reports Open Problems collection. The supplied proposed year is 2021; no individual proposer is named.
\item Supplied report citation: \emph{Partial Differential Equations (2022)}, Oberwolfach Reports, pp.\ 1859--1914, DOI \texttt{10.4171/owr/2021/35}. \url{https://doi.org/10.4171/owr/2021/35}. The record describes Oberwolfach Reports as documenting mathematical workshops and their open-problem discussions. The precise problem location within the report is not supplied.
\item Supplied literature source, labelled De Rosa--Tione: \url{https://arxiv.org/abs/2011.09922}. No title or theorem-level attribution is supplied in the record.
\item Supplied literature source, labelled Oberwolfach report: \url{https://ems.press/content/serial-article-files/46911}.
\item Literature assessment checked at \texttt{2026-08-22}: the supplied curation labels the problem partially solved and states that regularity is established for important graph classes and selected anisotropies, while the general anisotropic varifold regularity problem described by the record remains open. This is supplied dated curation derived from the OWR material, not an independently verified literature assessment; the status and cited claims require independent verification.
\item The registry labels the statement \texttt{corrected\_verified} and reports that the official report was checked directly and that definitions were restored while extraction artifacts and source-dependent wording were removed. This is a supplied verification claim, not verification performed for this draft; the supplied statement still leaves the qualifications listed above unspecified.
\end{itemize}
\end{document}
